\section{Collation: Three Narratives and One Message}

\subsection{The current custodial narrative}

The narrative a pilgrim meets today is fuller than either
nineteenth-century witness, and the National Shrine says plainly where
it comes from. Its closing paragraph is worth quoting, because it is
more careful than most retellings of it:

\begin{witnesstext}{National Shrine of Our Lady of Champion, ``Our
Story,'' closing note, read 2026-07-25}
Adele Brice, the visionary who witnessed Our Lady in Champion, never
learned to read or write and thus chronicled her experiences with the
Blessed Mother verbally. Because of this, the above apparition account
was developed based on a compilation of oral history from Adele and
sources close to her, historical documents, and second and third-hand
accounts found in the archives of the Shrine. While some details vary,
the account strives to do its best in determining the most common and
consistent information across the various historical
materials.\endnote{``Our Story,''
\url{https://championshrine.org/our-story/}, read 2026-07-25. The
shrine's separate historical timeline,
\url{https://championshrine.org/timeline/}, read the same day, supplies
the dated frame used in the table below: birth at Dion-le-Val, Belgium,
30 January 1831; a land purchase by Lambert and Marie Brice in the town
of Red River on 7 August 1855; the apparitions on 9 October 1859; a
second chapel in 1861 and the beginning of pilgrimages; a Third Order
Franciscan community, school, and convent in 1864; the fire of 8 October
1871; Adele's death on 5 July 1896. These are custodial statements, not
records this study verified.}
\end{witnesstext}

On that account, the first encounter is silent and happens while Adele
is walking a trail in the woods, her family later supposing she had
seen a soul in purgatory; the second, ``on what is believed to be
Sunday, October 9, 1859,'' occurs on the way to Mass with her sister
and a friend, at the same spot between two trees ``believed to be a
maple and hemlock''; the parish priest tells her, if it comes again, to
ask ``In God's name, who are you and what do you want of me?''; and the
third occurs on the road home, when the figure identifies herself as
the Queen of Heaven and gives the full message.

\subsection{The three narratives side by side}

\begin{collationtable}
Year & ``Twelve years ago'' from a text copyright-entered 1871: 1859 or
1858 & ``Ten or twelve years ago'' from an 1874 imprint: 1862--1864 &
1859 \\
Day and month & Not given; ``one summer day'' & Not given; three
successive Sundays & 9 October, ``believed to be'' a Sunday, with a
silent encounter some days earlier \\
Walk to Mass & Eleven miles & Two and a half leagues (English: seven
and a half miles) & Ten miles \\
Number and setting of encounters & Three: returning; going, a few weeks
later; returning that same afternoon & Three: all returning from Mass,
on three successive Sundays & Three: an errand in the woods; going to
Mass; returning from Mass \\
Companions & Not mentioned at the encounters & Present at the second
and third; ``a numerous escort, among whom were several men'' at the
third & Sister and a friend; they cannot see the figure and are told
to kneel \\
Priest consulted between encounters & No & No & Yes, after Mass, and he
supplies the form of the question \\
The figure & ``White and shining,'' face \emph{covered with a veil},
hands joined, in a luminous atmosphere, as in pictures of the
Immaculate Conception & ``A Lady of great beauty and majestic mien,''
suspended between two trees & Dazzling white, yellow sash, crown of
stars, long golden wavy hair, light so strong the face can hardly be
looked at \\
The two trees & Not mentioned & Mentioned, unidentified & Identified as
``believed to be a maple and hemlock'' \\
Who identifies the visitor & The seer recognizes ``the Virgin Mother''
& The seer addresses her as ``my good mother'' & The visitor: ``I am
the Queen of Heaven'' \\
The seer's opening words & ``Dear Mistress, what do you wish me to
do?'' & ``My good mother, what do you want of me?'' & ``In God's name,
who are you, and what do you want of me?'' \\
General confession, Communion for sinners & Absent & Absent & Present
\\
Warning of punishment & Absent & Absent & Present \\
Beatitude on the companions & Absent & Absent & Present \\
The vineyard rebuke & Absent & Absent & Present \\
Catechetical commission & Present, in full & Present, focused on First
Communion & Present, in Starr's wording \\
``Go and fear nothing, I will help you'' & Absent & Present & Present
\\
Annual commemoration & The Feast of the Assumption, 15 August & Twice a
year, ``at fixed epochs,'' unnamed & 9 October as the apparition
anniversary; 8 October for the fire; 15 August also kept \\
\end{collationtable}

\subsection{The stratigraphy of the reported words}

The most consequential result of the collation is textual rather than
narrative. Set the received message beside Starr's 1871 English and the
overlap is not thematic; it is verbal, and close to exact.

Starr: ``Gather around you the neglected children in this wild country,
running idle about the woods, and teach them what they should know for
salvation. \dots{} Teach them their catechism, how to sign themselves
with the sign of the cross, and how to approach the sacraments. This is
what I wish you to do.'' And Adele's objection: ``but how shall I teach
them, who know so little myself?''

The received text: ``Gather the children in this wild country and teach
them what they should know for salvation. \dots{} Teach them their
catechism, how to sign themselves with the sign of the Cross, and how
to approach the sacraments; that is what I wish you to do.'' And the
objection: ``But how shall I teach them who know so little myself?''

The received closing clause, ``Go and fear nothing, I will help you,''
is likewise Pernin's English, ``Go, and fear nothing; I will help
you.'' And the received narrative's ``wrapped as it were in a luminous
atmosphere'' reproduces Starr's distinctive phrase for the vision's
disappearance.

\begin{stratatable}
Catechetical commission (``gather the children\dots{} what they should
know for salvation''; catechism, sign of the cross, approach to the
sacraments; the seer's objection) & Starr, copyright-entered 1871 &
Received in Starr's English almost word for word; the earliest located
form of the message and its stable core \\
First Communion and the children who ``receive Him without knowing what
they do'' & Pernin, 1874 & Not in Starr; not in the received text
either, which has dropped it \\
``Go and fear nothing, I will help you'' & Pernin, 1874 & Received
verbatim; the one received clause with an 1870s witness that Starr
lacks \\
``Luminous atmosphere'' & Starr, 1871 & A narrative phrase, not a
spoken word, carried into the received narration \\
Self-identification as Queen of Heaven & Not located before the
received custodial narrative & Absent from both nineteenth-century
witnesses, in each of which the identification is made by the seer \\
General confession; Communion offered for the conversion of sinners &
Not located before the received custodial narrative & Absent from both
\\
``If they do not convert and do penance, my Son will be obliged to
punish them'' & Not located before the received custodial narrative &
Absent from both; also elided from the 2010 decree's own quotation \\
``Blessed are they that believe without seeing'' & Not located before
the received custodial narrative & Presupposes the companions'
question, which neither witness reports \\
``What are you doing here in idleness while your companions are working
in the vineyard of my Son?'' & Not located before the received
custodial narrative & Absent from both. Starr's only ``idle'' is the
children ``running idle about the woods'' \\
Crown of stars, yellow sash, golden hair & Not located before the
received custodial narrative & Starr's figure is explicitly veiled \\
Maple and hemlock & Not located before the received custodial narrative
& Pernin has two unidentified trees; the shrine says ``believed to be''
\\
\end{stratatable}

Two cautions belong immediately beside that table. First, ``not located
before the received custodial narrative'' means exactly that: this
study did not reach Sister M.~Dominica's shrine history of 1955, the
shrine's archival files, or any intermediate printing, and the elements
in question may well be attested in one of them at a date this study
cannot see. The claim is about the reach of the searches described in
Appendix~A, not about the history of the tradition. Second, the 2010
decree's quotation of the message includes the Queen of Heaven, the
general confession, and the Communion offered for sinners; it is the
punishment clause, the beatitude, and the vineyard rebuke that its
ellipses pass over. \S7.4 examines that quotation as a quotation.

\begin{judgmentbox}{Project synthesis --- the received message is a
composite with a datable core}
This study judges, on the verbal evidence set out above, that the
message now received at Champion is a composite English text whose
catechetical core reproduces the wording Eliza Allen Starr printed in
1871 and whose closing clause reproduces the wording Peter Pernin
printed in 1874, with a further stratum --- self-identification,
general confession, punishment, beatitude, vineyard --- that this study
could not trace to any nineteenth-century witness it reached. The
degree of verbal identity in the core is too high to be coincidence:
two independent English renderings of a Walloon conversation, made by
different people from different informants, do not converge on
``teach them what they should know for salvation'' and ``how to sign
themselves with the sign of the cross'' and ``that is what I wish you
to do.'' Somewhere in the chain, the received English descends from
printed English.\jsep
The strongest counterargument, stated at its strongest: Adele lived
until 1896 and told the story for thirty-seven years after Starr
printed it. A fixed English formula could have entered \emph{her own}
retelling --- read to her, or repeated to her by the sisters --- and
been transmitted from her lips as the received wording, in which case
the verbal identity proves textual continuity without proving that the
later strata are late. On that account the fuller message could be
material Starr and Pernin compressed away and Adele went on telling to
people who wrote it down after 1874. This is a serious reply and this
study cannot refute it from the sources it reached.\jsep
What would change the judgment: any dated witness earlier than the
1955 shrine history --- a letter, a parish or diocesan record, a
newspaper notice, a printed pamphlet, a chapel inscription --- carrying
the self-identification as Queen of Heaven, the command of general
confession, or the punishment clause. A single such witness would move
those elements out of the ``not located'' row and into the documented
tradition, and this study would record that it had been wrong about
their date.
\end{judgmentbox}

\subsection{The calendar of the cult, and where the anniversary moved}

One further result falls out of the collation, and it has gone
unremarked in every account of Champion this study saw. The date on
which the apparition is commemorated has moved.

Starr, the earliest witness, is unambiguous: ``Every year, on the Feast
of the Assumption of Our Lady, a religious procession celebrates the
appearance of the Blessed Virgin to Sister Adele.'' The commemoration
is annual, it is on 15 August, and its stated object is the apparition
itself. Pernin, three years later, records two solemn processions a
year ``at fixed epochs'' and does not name either date. The shrine's
own timeline places the beginning of ``devotions on the 15th of
August'' in 1871, in the entry for the fire --- which cannot be right
if Starr was describing an established annual custom in a volume
copyright-entered that same year.\endnote{``Shrine Historical
Timeline,'' read 2026-07-25, entry for 1871: ``Pilgrimages increase and
devotions on the 15th of August begin.'' Compare Starr, \work{Patron
Saints}, p.~301. This study does not attempt to date the origin of the
August observance; it records that the earliest witness located
presents it as already annual.}

Today the shrine keeps 8 October for the ``Miracle of the Great Fire''
with an all-night vigil, 9 October as the solemnity of Our Lady of
Champion and the apparition anniversary, 15 August for the Assumption,
and the first Saturday of May for the Walk to Mary. Of these, 9 October
is the one the shrine itself hedges: it is ``what is believed to be''
the day, and ``the date historians believe marks the anniversary of
Mary's last appearance.''\endnote{``Our Story'' and ``Solemnity of Our
Lady of Champion,'' both read 2026-07-25. The October pairing is
explicit on the solemnity page, which sets the vigil of 8 October
against the solemnity of 9 October.}

\begin{judgmentbox}{Project synthesis --- the October date is anchored
to the fire, not to the earliest cult}
This study judges that the earliest attested commemoration of the
Champion apparition was kept on 15 August, and that the October
anniversary is a later fixture whose gravitational centre is the fire
of 8 October 1871 rather than any early memory of the day of the
events. The reasons are that the only nineteenth-century witness to
name a date names 15 August and attaches it to the apparition; that
neither nineteenth-century witness gives any day or month for the
events at all; that the 2010 decree fixes the month of October and
conspicuously does not fix a day; and that the received day, 9 October,
sits immediately after the fire anniversary and is presented by the
shrine itself with ``believed to be.''\jsep
The strongest counterargument is that the two facts are independent:
the settlement might always have known the events fell in early
October, and the fire's falling on 8 October in the twelfth year
afterwards --- which the shrine describes as ``almost twelve years to
the date'' --- would then be the coincidence that fixed public
attention on the date rather than the cause of it. A community can
easily keep a summer procession on a great Marian feast for
convenience while remembering an autumn anniversary privately. That is
entirely possible, and it is why this study's claim is about
attestation rather than about origin.\jsep
What would change the judgment is any pre-1871 source giving a day or
month for the events, or any pre-1955 source giving 9 October as the
apparition's day independently of the fire's anniversary.
\end{judgmentbox}

\subsection{What none of this touches}

The stratigraphy above is a statement about texts. It is not a
statement about whether the events occurred, and it is not a challenge
to the 2010 decree, which judges ``the substance'' of the events and
does not adopt any particular wording as authentic. It is also not a
reason to stop using the received message in prayer. A devotional text
whose earliest located form is a catechetical commission and whose
later form calls for conversion, penance, and confession is not
teaching anything the Church does not teach in its own name. What the
stratigraphy forbids is one specific move: treating any sentence of the
received message as a transcript of what was said in a Wisconsin wood
in 1859, and then reasoning from the exact words as if they had the
authority of Scripture. That move is forbidden anyway, for every
private revelation, on grounds that have nothing to do with textual
criticism.
